\section{Conclusiones y Comentarios}

En términos generales, el modelo sísmico desarrollado en ETABS es consistente con las bases de 
cálculo del prediseño y cumple con los requisitos establecidos en la NCh433.

En cuanto a las deformaciones, ambas direcciones cumplen con el límite admisible de 
deriva de entrepiso $\theta \leq 0{,}002$. La dirección X presenta la mayor demanda con 
$\theta_{x,\text{máx}} = 0{,}001468$ en el entrepiso 14--15, y la dirección Y alcanza un máximo 
de $\theta_{y,\text{máx}} = 0{,}000909$ en el entrepiso 11--12. Los desplazamientos de techo 
son 67,18~mm y 43,19~mm en X e Y respectivamente, valores que se consideran como un comportamiento 
lateral adecuado para un edificio de esta altura.

Respecto de la rigidez, la verificación global mediante la razón $H/T$ muestra que la 
dirección Y (64,22~m/s) se encuentra dentro del rango esperado de 50 a 70~m/s, mientras que 
la dirección X (48,05~m/s) queda levemente por debajo, lo que es coherente con su mayor 
período fundamental. Por piso, la rigidez lateral disminuye de forma progresiva en altura, 
concentrándose la mayor variación entre los pisos 1 y 2 ($K_{x,2}/K_{x,1} \approx 0{,}66$ 
y $K_{y,2}/K_{y,1} \approx 0{,}69$), sin que se identifique un piso blando.

En lo que respecta al acoplamiento, el edificio no cumple completamente la 
verificación de desacoplamiento modal. La relación $T_{\text{tors}}/T_x = 1{,}10$ no supera 
el límite de 1,2, lo que indica que el modo torsional y el traslacional en X están 
relativamente cercanos entre sí, con períodos de 1,251~s y 1,140~s, respectivamente. Esto 
sugiere que la respuesta sísmica en esa dirección podría verse influenciada por efectos 
torsionales. Las otras dos relaciones sí cumplen: $T_{\text{tors}}/T_y = 1{,}47$ y 
$T_x/T_y = 1{,}34$.

En cuanto a los cortes basales y pesos sísmicos, el peso sísmico obtenido del modelo 
es $P = 7.660{,}05$~tonf, lo que representa una diferencia del 12,6\% respecto del valor de 
prediseño de 8.766,65~tonf, diferencia que se encuentra dentro de los márgenes esperados para 
esta etapa del análisis. Los cortes basales obtenidos en el nivel 1 son 746,41~tonf en X y 
753,54~tonf en Y, frente al corte objetivo de 766,05~tonf. Las diferencias son de un 2,6\% 
y 1,6\%, respectivamente, lo que indica que el modelo es consistente con las bases de cálculo.


Finalmente, respecto del período principal de traslación, el modelo entrega 
$T_X = 1{,}141$~s en dirección X y $T_Y = 0{,}854$~s en dirección Y. La diferencia con el 
período de prediseño ($T = 1{,}05$~s) es de un 8,7\% en la dirección X, lo cual es razonable 
dado que la expresión empírica utilizada en el prediseño no distingue dirección de análisis 
ni incorpora la rigidez real de los muros modelados.